\chapter{Features}
\label{ch:features}

This chapter describes every feature of the app: what the user sees, what happens, and where it is
implemented. Each feature ends with a summary of its screens (under \code{src/app}), the client functions
that carry it, and the server rules that enforce it. The data flows behind the larger features are traced
in \cref{ch:flows}; this chapter refers to them rather than repeating them. \Cref{fig:screen-map} shows
how the signed-in screens connect, and \cref{app:routes} lists every route.

\begin{figure}[ht]
  \centering
  % Screen map: how the signed-in screens connect, with the action on each transition (Figure: screen map).
\begin{tikzpicture}[
    tab/.style={box, accentbox, font=\sffamily\scriptsize\bfseries, text width=20mm, minimum height=9mm, align=center},
    scr/.style={box, font=\sffamily\scriptsize, text width=20mm, minimum height=9mm, align=center},
    act/.style={font=\sffamily\tiny, text=muted, fill=white, inner sep=1pt, align=center},
  ]
  \node[tab] (golive) at (1.2,0) {Go live tab};
  \node[tab] (nearby) at (4.6,0) {Nearby tab};
  \node[tab] (map) at (8.5,0) {Map tab};
  \node[tab] (profile) at (13.45,0) {Profile tab};

  \node[scr] (sheet) at (1.2,-2.0) {Go-live sheet\\\texttt{go-available}};
  \node[scr] (person) at (4.6,-2.0) {Person sheet\\\texttt{person/[id]}};
  \node[scr] (waiting) at (4.6,-4.0) {Waiting\\\texttt{waiting}};
  \node[scr] (incoming) at (1.2,-4.0) {Incoming request\\\texttt{incoming}};
  \node[scr, accentbox] (chat) at (2.9,-6.0) {Chat\\\texttt{chat}};
  \node[scr] (newev) at (7.2,-2.0) {Event form\\\texttt{event/new}};
  \node[scr] (event) at (9.75,-2.0) {Event sheet\\\texttt{event/[id]}};
  \node[scr, accentbox] (group) at (8.5,-4.0) {Group chat\\\texttt{group/[id]}};
  \node[scr, draw=alert, fill=alertsoft] (report) at (8.5,-6.0) {Report\\\texttt{report/[id]}};
  \node[scr] (edit) at (12.2,-2.0) {Edit profile\\\texttt{edit-profile}};
  \node[scr] (settings) at (14.7,-2.0) {Settings\\\texttt{settings}};
  \node[scr, text width=24mm, align=center] (rows) at (14.5,-4.1) {\texttt{password}, \texttt{blocked},\\\texttt{tutorial}; sign out;\\delete account};

  \draw[flow] (golive) -- node[act] {tap YAP,\\or Edit note} (sheet);
  \draw[flow] (nearby) -- node[act] {tap a card\\or a prompt} (person);
  \draw[flow] (person) -- node[act] {Say hi} (waiting);
  \draw[flow] (waiting) |- node[act, pos=0.25] {accepted} (chat.east);
  \draw[flow] (incoming) |- node[act, pos=0.25] {Accept} (chat.west);
  \draw[flow] (map) -- node[act] {long press, or\\Create event here} (newev);
  \draw[flow] (map) -- node[act] {tap a pin} (event);
  \draw[flow] (newev) -- node[act] {Post event} (group);
  \draw[flow] (event) -- node[act] {Join} (group);
  \draw[flow] (profile) -- node[act] {Edit profile} (edit);
  \draw[flow] (profile) -- node[act] {gear} (settings);
  \draw[flow] (settings) -- (settings |- rows.north);
  \draw[softflow] (group) -- node[act] {long press\\a message} (report);

  \node[act, anchor=south, text width=22mm, fill=none] at ($(incoming.north)+(0,0.05)$) {opens over any screen\\when a request arrives};
  \node[act, anchor=north, text width=46mm, fill=none] at ($(report.south)+(0,-0.05)$) {also from the safety menu on the person sheet,\\incoming request, chat and event sheet};
  \node[act, anchor=north, text width=40mm, fill=none] at ($(chat.south)+(0,-0.05)$) {reopened from the bar on the Go live and Nearby tabs};
\end{tikzpicture}

  \caption{Screen map of the signed-in app. Green boxes along the top are the four tabs; arrows carry the
    action that opens the next screen. The incoming request and the chat can open over any screen.}
  \label{fig:screen-map}
\end{figure}

\section{Welcome and account creation}
\label{sec:feature-signup}

The welcome screen shows the name \ui{Yap}, the line \ui{The app that gets you off the app}, and two
buttons: \ui{Create account} and \ui{Sign in}. A note states that continuing means agreeing to the
Community Guidelines and Privacy Policy. \ui{Create account} asks for an email address, a password and its
confirmation. Above its button, a line says that creating an account means agreeing to the Terms of
Service and the Community Guidelines and confirming the Privacy Policy has been read; each name is a link
(the terms and privacy policy open in the browser, the guidelines in the app). The button then opens the
code screen (\cref{sec:flow-account}). The code screen shows eight boxes,
fills them as digits are typed or pasted, submits on the eighth digit, and offers \ui{Resend code} after
a 60-second countdown.

\implemented
  {\code{(auth)/welcome}, \code{(auth)/sign-up}, \code{(auth)/verify}}
  {\code{signUp}, \code{confirmSignUp}; \code{api.resendSignUpCode}; \code{passwordProblem}}
  {Supabase Auth: email confirmation on, minimum password length 8, 8-digit code in the confirmation template}

\section{Signing in and passwords}

\ui{Sign in} takes an email and a password. An unconfirmed account is sent a new code and taken to the
code screen. \ui{Forgot password?} carries the typed email to the reset screen, which emails a reset code;
the code screen then leads to \ui{Choose a new password}, and saving signs the user in
(\cref{sec:flow-account}). A note on the sign-in screen tells users who signed up with an emailed code
alone, before passwords existed, to set a password this way. Signed-in users change their password from
\ui{Settings}, \ui{Change password}; a success alert confirms the change. A new password must have
at least eight characters and match its confirmation before anything is sent.

\implemented
  {\code{(auth)/sign-in}, \code{(auth)/forgot-password}, \code{(auth)/verify}, \code{(auth)/new-password}, \code{(app)/password}}
  {\code{signInWithPassword}, \code{finishPasswordReset}; \code{api.requestPasswordReset}, \code{api.setPassword}}
  {Supabase Auth: password sign-in, recovery codes, at most one code per address per minute on the hosted project (\cref{tab:dashboard})}

\section{Profile setup}
\label{sec:feature-setup}

A new account first sees a greeting and must tick \ui{I'm 18 or older} before continuing. Six numbered
steps follow, one question per screen, each with a progress bar and its number:

\begin{enumerate}
  \item first and last name; the screen notes that others see only the first name and last initial;
  \item gender (woman, man, non-binary or prefer not to say) and a switch to show it on the profile, which
    is disabled for \ui{Prefer not to say};
  \item an introduction of up to 300 characters;
  \item three to five interests from a list of 32;
  \item a profile picture, chosen from the library or taken with the front camera, cropped square;
  \item four to six further photos, picked several at a time, each of which can be replaced or removed.
\end{enumerate}

\ui{Finish} saves the profile (\cref{sec:flow-profile}). The greeting also offers
\ui{Use a different account}, which signs out, for someone who signed up with the wrong address.

\implemented
  {\code{profile-setup/index}, \code{name}, \code{gender}, \code{about}, \code{interests}, \code{photo}, \code{photos}}
  {\code{ProfileDraftProvider}, \code{saveProfile}, \code{api.saveProfile}; \code{PhotoGrid}, \code{SetupStep}}
  {\code{profiles} check constraints (lengths; names and intro with something visible and no control characters, direction overrides or stacked accents; 3 to 5 different interests of 1 to 30 characters; own-folder photo paths; at most 6 extra photos); insert policy requires \code{is_adult}}

\section{Profile tab and editing}

The \ui{Profile} tab shows the user's profile exactly as other people see it, built by
\code{asSeenByOthers} from the stored profile: first name and last initial, gender only if shown, the live
note while the user is live, the further photos, the introduction and the interests. A gear opens
\ui{Settings} and \ui{Edit profile} opens a single-page form with every setup field except the 18+
confirmation, which keeps its original answer. The form lists anything still missing (\ui{Add at least 4
more photos, \dots{} to continue}) and saves through the same action as setup.

\implemented
  {\code{(app)/(tabs)/profile}, \code{(app)/edit-profile}}
  {\code{ProfileForm}, \code{PersonDetails}, \code{asSeenByOthers}, \code{saveProfile}}
  {Column grants: a client may update only the editable profile columns}

\section{Going live}
\label{sec:feature-live}

The \ui{Go live} tab holds one large round button. Green, it reads \ui{YAP}; tapping it opens the go-live
sheet. There the user picks 30, 60 or 120 minutes (60 is preselected) and may add a note of up to 60
characters, with a counter. A notice explains that the location is used only to find people in the same
place while the user is available, that nobody sees it, and that it is deleted on stopping. If location
permission is refused, the notice changes to a warning and the button to \ui{Open Settings}.

While the user is live the button is red and reads \ui{LIVE}, the minutes left and \ui{Tap to get off};
its pulsing ring speeds up. Below it the note appears with \ui{Edit} (or \ui{Add} when empty), which opens
the same sheet in edit mode; saving the note updates the screen at once and restores the old note if the
server refuses. A caption reminds the user to keep Yap open to receive requests, and the screen stays on.
When the time runs out, the user goes offline and a dismissible \ui{Your time's up} notice appears.

\implemented
  {\code{(app)/(tabs)/index}, \code{(app)/go-available}}
  {\code{HeroButton}, \code{DurationChips}, \code{useForegroundLocation}, \code{goGreen}, \code{updateNote}, \code{goOffline}; keep-awake and the end timer in \code{(app)/_layout}}
  {\code{go_green} (durations, note length, relocation limits), \code{update_note}, \code{go_offline}; \code{expire_stale} ends sessions on time}

\section{Nearby list and prompts}
\label{sec:feature-nearby}

The \ui{Nearby} tab requires the user to be live, because only people who can see each other appear to
each other; otherwise it shows the Yap cat and a \ui{Go live} button. While live, it shows \ui{You're live,
47 min left} and a card for each nearby person: photo, first name and last initial, a \ui{here now} badge,
the note, two lines of introduction and up to three interests. Placeholder cards pulse while the first
list loads; an empty list shows \ui{Nobody else is green here yet}. Above the list, a bar shows a chat in
progress and another the request the user is waiting on, with its countdown, or \ui{Not this time} once
it has ended. While live, a prompt slides in at the top of any screen when someone new is nearby, with
\ui{Say hi} or \ui{See them} (\cref{sec:flow-live}).

\implemented
  {\code{(app)/(tabs)/nearby}; prompts in \code{(app)/_layout}}
  {\code{useNearbyPolling}, \code{useNearby}, \code{refreshNearby}, \code{useNearbyAlerts}, \code{NearbyToast}, \code{PersonCard}}
  {\code{nearby}: live, open, active users in the same cell of about 100\,m or a touching one, no blocks, at most 50, cards only, 12 calls a minute}

\section{Person sheet and saying hi}

Opening a person shows a sheet with their full profile: large photo, name, gender if shown, note, further
photos (signed only now), introduction and interests, with the safety menu in the corner. The sheet looks
the person up in the shared nearby list, the current chat partner or the sender of the request on screen;
if they have gone, it says \ui{They're not live any more}. The footer offers \ui{Say hi}, unless the user is
already chatting (with this person or anyone else) or has a request waiting, in which case it explains
why. A sent request replaces the sheet with the waiting screen.

\implemented
  {\code{(app)/person/[id]}}
  {\code{PersonDetails}, \code{usePersonPhotos}, \code{sayHi}, \code{api.sendRequest}}
  {\code{send_request}: proximity, blocks, one request waiting, one chat, 5 requests per 10 minutes}

\section{Waiting for an answer}

The waiting screen shows the recipient's photo with the Yap cat beside it, \ui{Waiting for Sam\dots},
\ui{They have 5 minutes to answer} and a countdown, with a reassurance that \ui{Not this time} carries no
score. \ui{Back to nearby list} leaves it open in the background, where the bar on the Nearby tab tracks
it. \ui{Cancel request} withdraws it on the server. When the answer is no, or time runs out, the screen
changes to \ui{Not this time} (\cref{sec:flow-sayhi}).

\implemented
  {\code{(app)/waiting}}
  {\code{watchAnswer}, \code{api.waitForAnswer}, \code{cancelRequest}, \code{expireOutgoing}, \code{clearOutgoing}}
  {\code{cancel_request}; \code{expire_stale} expires the request after five minutes or when the sender stops being live}

\section{Answering a request}

An incoming request takes over the screen wherever the user is, with a haptic. It shows \ui{Someone wants
to chat}, the sender's full profile and the time left to answer, and offers \ui{Not now} and \ui{Accept}
with equal size and weight. The safety menu is in the corner. The screen cannot be swiped away; it closes
when the user answers, when the time runs out, or when the user blocks the sender.

\implemented
  {\code{(app)/incoming}; opened by \code{(app)/_layout}}
  {\code{listenForRequests}, \code{onIncoming}, \code{respondToIncoming}, \code{expireIncoming}}
  {\code{respond}: one chat per user, the sender still live; \code{request_person}: card hidden after a block}

\section{One-to-one chat}

An accepted request opens the chat for both users. The header shows the partner's photo and name with
\ui{Your yap partner --- go say hi}; tapping it opens their profile. A new chat shows the happy Yap cat
above the first message, and in a user's first chats a safety tip, shown until it has been seen three
times, advises meeting in the public place both are in. Messages up to 1000 characters are sent with the arrow button; the user's own messages
are green and on the right. \ui{End} asks whether the two met. If the partner ends the chat, the
composer is replaced by \ui{Sam ended the chat} and \ui{End} becomes \ui{Close}. The chat survives leaving
the screen, and a bar on the \ui{Go live} and \ui{Nearby} tabs reopens it (\cref{sec:flow-sayhi}).

\implemented
  {\code{(app)/chat}}
  {\code{openChat}, \code{api.subscribeChat}, \code{sendMessage}, \code{endChat}, \code{markChatEnded}; \code{ChatThread}, \code{ActiveChatBar}}
  {\code{send_chat_message} (length, membership, blocks, 30 per minute), \code{end_chat}; Realtime delivery limited by the \code{messages} policy}

\section{Map and events}
\label{sec:feature-events}

The \ui{Map} tab shows events around the user. A banner explains that pins show events and never people,
and holds the distance slider (1 to 50\,km). A button above \ui{Create event here} recentres the map on
the user. Pins show the event's category icon and title; joined events are green. The event form, the
event sheet and the group chat behave as described in \cref{sec:flow-events}. \Cref{tab:event-rules}
summarises who may do what.

\begin{table}[ht]
  \centering
  \caption{Event rules.}
  \label{tab:event-rules}
  \tablefont
  \begin{tabularx}{\textwidth}{@{}L{34mm}Y@{}}
    \toprule
    \textsf{\textbf{Action}} & \textsf{\textbf{Rule}} \\
    \midrule
    See an event & Within the chosen radius, at most 50\,km, of the position the server holds for the user; until three hours after its start; host active and not blocked either way; the host's photo only after joining \\
    Create an event & Active adult profile; one of the twelve categories; start from five minutes ago to seven days ahead; at most five per host in 24 hours \\
    Join an event & Active adult profile; within 5\,km of the position the server holds; not blocked by or blocking the host \\
    Read and write its chat & Members only; 1 to 1000 characters; at most 30 messages a minute \\
    Leave & A guest's membership ends; the app offers no \ui{Leave} to the host \\
    Delete & Host only, from the event sheet, after a confirmation; members and the group chat go with the event \\
    \bottomrule
  \end{tabularx}
\end{table}

\implemented
  {\code{(app)/(tabs)/map}, \code{(app)/event/new}, \code{(app)/event/[id]}, \code{(app)/group/[id]}}
  {\code{loadEvents}, \code{setEventRadius}, \code{createEvent}, \code{joinEvent}, \code{leaveEvent}, \code{deleteEvent}, \code{enterEventChat}, \code{sendGroupMessage}; \code{EventIcon}, \code{eventCategory}}
  {\code{set_map_position}, \code{events_near}, \code{create_event}, \code{join_event}, \code{leave_event}, \code{delete_event}, \code{send_event_message}}

\section{Blocking and reporting}
\label{sec:feature-safety}

Every surface that shows another person (the person sheet, the incoming request, the chat and the event
sheet) has a menu with \ui{Report} and \ui{Block}; in a group chat, a long press on a message offers the
same. Blocking asks for confirmation, then takes effect on the server first: the block is recorded and any
request or chat between the two is ended. The app then drops the person from the nearby list, the map and
any request or chat on screen. \ui{Report} offers six reasons, an optional comment of up to 500 characters
and \ui{Also block}, which is ticked by default. The form says that every report is reviewed and that the
reported person will not know who reported them. \ui{Settings}, \ui{Blocked people} lists blocked
users by first name and last initial, with \ui{Unblock}. The server semantics are in
\cref{sec:blocking-reporting}. Separately, the server refuses an introduction, interest, live note, event
or message that contains a slur or abusive word from its list, and the app explains why
(\cref{sec:word-filter}); the owner reviews reports and can suspend an account (\cref{sec:moderation}).
A hidden or suspended user sees \ui{Your account is on hold} instead of the app, with \ui{Email support},
\ui{Use a different account} and \ui{Delete account}, which opens the same delete screen as Settings; the
screen appears at launch, and as soon as the server refuses an
action with \code{account restricted}, and is checked again at each catch-up and whenever the delete screen
opens, so a restored account comes back by itself and the delete screen says what is kept.

\implemented
  {\code{account-restricted} (a root stack)}
  {\code{restricted}, \code{checkRestricted}; \code{api.fetchMyAccount}, \code{api.onAccountRestricted}}
  {\code{private.require_active} in every function a user acts through}

\implemented
  {\code{SafetyMenu}, \code{(app)/report/[id]}, \code{(app)/blocked}, long press in \code{(app)/group/[id]}}
  {\code{blockUser}, \code{unblockUser}, \code{reportUser}; \code{api.block}, \code{api.unblock}, \code{api.blockedList}, \code{api.report}}
  {\code{block_user}, \code{blocked_list}, \code{report_user}; delete policy on \code{blocks}; block checks in every function that involves two users; the word-filter triggers}

\section{Settings, signing out and deleting the account}

\ui{Settings} groups its rows like the phone's own settings. \textsf{Account} holds \ui{Change password},
\ui{Blocked people} (with the count) and \ui{Notification settings}, which opens the app's page in the
phone's settings. \textsf{Help} holds \ui{View tutorial}. \textsf{Safety and legal} holds
\ui{Community guidelines} and \ui{Acknowledgements}, which open screens in the app, and \ui{Privacy
policy}, \ui{Terms of service} and \ui{Contact support}, which open the addresses in
\code{src/config/links.ts} (\code{https://yap.samarxkx.dev/privacy}, \code{/terms} and an email to
\code{support@samarxkx.dev}). The owner publishes the pages from \code{docs/legal/} at those addresses
before the first store release (\cref{sec:release}). Below are \ui{Sign out}, \ui{Delete account} and
the app version.

\ui{Community guidelines} lists the rules everyone agrees to (adults only, be respectful, keep it clean, be
yourself, nothing illegal or dangerous, meet safely) and what happens when someone breaks them: reports
are reviewed within 24 hours, content that breaks the rules is removed, people who break them are
suspended or removed and cannot sign up again with the same email, a profile reported by three people who
really interacted with the person on Yap (or have used it for a few days) is hidden straight away, and slurs and abusive
words are refused when posted. \ui{Acknowledgements} shows Yap's copyright and every open-source package
the app ships, each with its licence text, from the list \code{npm run licenses} generates.

Signing out requires a connection; offline, it explains why and does nothing, because ending the chat and
the live session must reach the server or the other person would be left waiting. Online, it ends any chat,
goes offline, then signs out, and the Auth listener resets the app. \ui{Delete account} opens its own
screen, which lists what is deleted and keeps \ui{Delete my account} disabled until the user types their
first name as it is on their profile, ignoring capitals and spaces around it; deleting then follows
\cref{sec:flow-delete}. A hidden or suspended user reaches the same screen from \ui{Your account is on
hold}. This in-app deletion is what Apple requires (guideline 5.1.1(v)); for Google Play,
the published \code{delete-account} page also tells people without the app how to ask for deletion.

\implemented
  {\code{(app)/settings}, \code{(app)/password}, \code{(app)/blocked}, \code{(app)/delete-account}, \code{guidelines}, \code{acknowledgements}}
  {\code{signOut}, \code{deleteAccount}; \code{api.deleteAccount}; \code{LINKS}; \code{third-party-licenses.json}}
  {\code{delete_account}; Storage delete policy on the user's own folder}

\section{Tutorial}

A new account sees a five-page tutorial once, right after saving its first profile: going live, seeing who
is around, saying hi and meeting, making plans on the map, and staying in control. It counts as seen as
soon as it opens, so leaving it by any route never brings it back. \ui{Settings}, \ui{View tutorial} shows it
again.

\implemented
  {\code{(app)/tutorial}; opened by \code{(app)/_layout}}
  {\code{tutorialDue}, \code{markTutorialSeen}}
  {None}

\section{Offline and error states}
\label{sec:feature-offline}

A banner, \ui{No internet connection. Trying to reconnect}, appears at the top of every screen whenever
the phone reports no connection. If a signed-in user's profile cannot be loaded, the app shows \ui{Can't
reach Yap} with \ui{Try again} instead of sending them to profile setup, and recovers by itself when a
later catch-up succeeds. Every refused or failed action shows a sentence, never database text
(\cref{app:refusals}). \Cref{sec:offline} describes the offline behaviour of each action.

\implemented
  {\code{OfflineBanner}, \code{CantLoad} (both in the root layout)}
  {\code{retryLoad}, \code{catchUp}; \code{SERVER_REFUSALS}, \code{friendlyRpcError}, \code{friendlyAuthError}}
  {Fixed refusal messages in every database function}

\section{Visual design and accessibility}

The theme in \code{src/theme/tokens.ts} defines a light and a dark palette, chosen to follow the phone's
setting, in which every text and background pair meets WCAG AA contrast (4.5:1 for text, 3:1 for
interface parts). Spacing, radii, sizes, type styles and two elevations come from shared scales. Text
uses the system fonts, so the phone's text-size setting applies; the largest display words cap their
scaling. Tappable surfaces share one press animation that shrinks and darkens them slightly; loading lists
show pulsing placeholders instead of spinners; tab changes and key actions play light haptics.

The Yap cat appears on empty, waiting and error screens in five poses (dizzy, waiting, searching,
sleepy and happy), each with a slow loop of movement. The cat, the pulsing button ring and the press
animation all stop moving when the phone's reduced-motion setting is on, and the cat is hidden from screen
readers because the text beside it always says what it means. Interactive elements carry accessibility
roles, labels and states, the code boxes announce how many digits are entered, and minimum touch targets
are 44 points.

\implemented
  {All}
  {\code{useTheme}, \code{usePressFeedback}, \code{Skeleton}, \code{Mascot}, \code{GlassBackdrop}}
  {None}

\section{Feature built but switched off}
\label{sec:feature-flags}

One feature is present in the code and disabled by a flag in \code{src/config/features.ts}.

\paragraph{Sign-in with Apple and Google} When enabled, the welcome screen shows \ui{Continue with Apple}
(on iOS) and \ui{Continue with Google}. Neither provider is configured in Supabase Auth, and the store's
handler only shows \ui{Not available yet}.

\paragraph{Selfie and face verification} This release has none. The screen, the placeholder matcher, the
phone-side verified state and the \ui{Verified} badge were removed from the app. The code is parked on the
\code{future/face-verification} branch and tracked as issue 17. The database column \code{profiles.verified}
remains, but the app ignores it and no screen shows or asks for verification.
